\section{Conclusion}
\label{sec:conclusion}

A valuation framework has been developed for expiry-hour electricity options in a
market that supplies a sparse monthly forward strip and no option prices. Its
organising principle is a separation between what such a market identifies and
what it does not. An equality-constrained hourly curve carries the identified
information, a regime-switching residual with residual-demand-driven transition
probabilities carries the assumed distributional structure, and a centring term
ensures the two meet without the second contaminating the first:
$\E^{\Q}[P_t]=F(t)$ holds for every admissible parameter vector, as an algebraic
identity rather than a calibration target.

Applying the framework to Turkish data yields three findings. The dominant source
of uncertainty in the reported option value is the forward level over the
unquoted nearest delivery month, which overwhelms every parameter of the residual
specification by an order of magnitude. The mean-reversion rate is the dominant
residual-layer input, and transferring it from an autoregression on the
untransformed price level rather than on deviations around the curve inflated
option values by a factor of four until diagnosed. And the forward strip itself is a poor predictor of what follows: across seven
valuation dates the mean error is $-18\%$ of the forward level, negative at five
dates and positive at two, with the much-discussed 2026 outcome lying inside the
interquartile range rather than outside it.

The last of these deserves emphasis. Forward-anchored valuation inherits the
predictive accuracy of the forward market entirely, and on this evidence that
accuracy is low and unstable in sign for reasons no pricing model can correct.
Combined with the finding that the nearest delivery month is never quoted, the
implication for practice is that a single option value in this market conveys
less than the value reported together with the range implied by the inputs no
price identifies. What that discipline looks like in practice is what this paper
has tried to demonstrate, including where it revealed our own errors.
